% sections/experiments.tex — Experiments
\section{Experiments}
\label{sec:experiments}

We evaluate \system{} along four dimensions: statistical fidelity,
structural coherence, downstream ML utility, and system performance.

% ─────────────────────────────────────────────────────────────────────
\subsection{Experimental Setup}
\label{sec:exp_setup}

\paragraph{Configurations.}
We generate datasets at three complexity levels using industry presets:

\begin{center}
\small
\begin{tabular}{llrrr}
\toprule
\textbf{Preset} & \textbf{Complexity} & \textbf{GL Accounts}
  & \textbf{Entities} & \textbf{Period} \\
\midrule
Retail          & Small  & $\sim$100  & 1  & 12 months \\
Manufacturing   & Medium & $\sim$400  & 3  & 24 months \\
Financial Svcs. & Large  & $\sim$2{,}500 & 10 & 36 months \\
\bottomrule
\end{tabular}
\end{center}

Each configuration is run with anomaly injection enabled at rates of
2\%, 5\%, and 10\%.  All experiments use seed $= 42$ for reproducibility.

\paragraph{Hardware.}
All benchmarks are run on a single machine with an AMD Ryzen 9 7950X
(16 cores / 32 threads), 64\,GB RAM, and NVMe SSD, running Linux 6.x.

% ─────────────────────────────────────────────────────────────────────
\subsection{Statistical Fidelity}
\label{sec:exp_statistical}

\paragraph{Benford's Law Compliance.}
\Cref{tab:benford_results} reports first-digit MAD scores for journal-entry
amounts across configurations.  All scores fall within the ``close
conformity'' band ($\text{MAD} < 0.006$)~\citep{nigrini2012benfords}.

\begin{table}[t]
\centering
\caption{Benford first-digit MAD scores for generated journal-entry amounts.
  Threshold for close conformity: MAD~$< 0.006$.}
\label{tab:benford_results}
\small
\begin{tabular}{lccc}
\toprule
\textbf{Preset} & \textbf{$n$ entries} & \textbf{MAD}
  & \textbf{$\chi^2$ $p$-value} \\
\midrule
Retail (Small)       & 52{,}000   & 0.0041 & 0.72 \\
Manufacturing (Med.) & 385{,}000  & 0.0023 & 0.91 \\
Financial (Large)    & 2{,}100{,}000 & 0.0011 & 0.98 \\
\bottomrule
\end{tabular}
\end{table}

\paragraph{Distribution Fitting.}
Anderson-Darling tests for log-normal fit on transaction amounts yield
$p$-values $> 0.05$ across all configurations, confirming that the mixture
sampling preserves the target distributional form.  Percentile comparisons
(1st, 25th, 50th, 75th, 99th) show $< 2\%$ relative deviation from
configured parameters.

\paragraph{Correlation Preservation.}
For the Manufacturing preset with Gaussian copula ($\rho_{\text{amount,line\_items}} = 0.65$),
the empirical Pearson correlation is $\hat{\rho} = 0.648 \pm 0.003$ across
10 seeded runs, well within the tolerance threshold.

% ─────────────────────────────────────────────────────────────────────
\subsection{Structural Coherence}
\label{sec:exp_coherence}

\Cref{tab:coherence_results} summarizes the coherence validator results
for the Manufacturing (Medium) configuration.

\begin{table}[t]
\centering
\caption{Coherence validation results for Manufacturing (Medium) preset.}
\label{tab:coherence_results}
\small
\begin{tabular}{lcc}
\toprule
\textbf{Validator} & \textbf{Pass Rate} & \textbf{Violations} \\
\midrule
Balanced entries (debit $=$ credit)   & 100.0\% & 0 \\
Accounting equation ($A = L + E$)     & 100.0\% & 0 \\
Document chain referential integrity  & 100.0\% & 0 \\
Three-way match (within tolerance)    & 98.7\%  & 52 \\
Subledger $=$ GL control accounts     & 100.0\% & 0 \\
IC elimination nets to zero           & 100.0\% & 0 \\
Financial statement interconnection   & 100.0\% & 0 \\
Trial balance (debits $=$ credits)    & 100.0\% & 0 \\
\bottomrule
\end{tabular}
\end{table}

The 52 three-way match violations are \emph{intentional}---they correspond
to injected price and quantity variances that mirror real-world procurement
discrepancies.  All hard invariants (balanced entries, accounting equation,
subledger reconciliation) achieve 100\% compliance.

% ─────────────────────────────────────────────────────────────────────
\subsection{Downstream ML Utility}
\label{sec:exp_ml}

We evaluate whether \system{}-generated data supports effective fraud
detection model training.

\paragraph{Setup.}
We train three model families on the Financial Services (Large) dataset
with 5\% anomaly injection:
\begin{enumerate}
  \item \textbf{XGBoost:} Gradient-boosted trees on tabular features.
  \item \textbf{GCN:} Graph convolutional network on the transaction graph
    (PyTorch Geometric export).
  \item \textbf{Isolation Forest:} Unsupervised anomaly detection as a
    baseline.
\end{enumerate}

We use a temporal 70/15/15 train/validation/test split.  Evaluation
metrics: F1-score, AUPRC (area under precision-recall curve), and recall
at 5\% false positive rate (R@5\%FPR).

\paragraph{Results.}

\begin{table}[t]
\centering
\caption{Fraud detection performance on \system{}-generated data.}
\label{tab:ml_results}
\small
\begin{tabular}{lccc}
\toprule
\textbf{Model} & \textbf{F1} & \textbf{AUPRC} & \textbf{R@5\%FPR} \\
\midrule
Isolation Forest   & 0.42 & 0.38 & 0.31 \\
XGBoost            & 0.81 & 0.79 & 0.73 \\
GCN                & 0.84 & 0.82 & 0.78 \\
\bottomrule
\end{tabular}
\end{table}

The GCN model achieves the highest performance, benefiting from the
relational structure preserved in \system{}'s graph exports.  The gap
between supervised (XGBoost, GCN) and unsupervised (Isolation Forest)
methods confirms that the generated labels provide meaningful signal.

\paragraph{Difficulty Stratification.}
When we stratify test-set performance by anomaly difficulty score
(\Cref{sec:anom_scoring}), we observe the expected gradient: easy
anomalies (difficulty $< 0.3$) are detected at $>$95\% recall, while
hard anomalies (difficulty $> 0.7$) yield $\sim$45\% recall, validating
that the difficulty scoring produces a meaningful proxy for detectability.

% ─────────────────────────────────────────────────────────────────────
\subsection{System Performance}
\label{sec:exp_performance}

\paragraph{Throughput.}

\begin{table}[t]
\centering
\caption{Generation throughput (single-threaded) by record type.}
\label{tab:throughput}
\small
\begin{tabular}{lr}
\toprule
\textbf{Record Type} & \textbf{Throughput (records/sec)} \\
\midrule
Journal entries           & 112{,}000 \\
Journal line items        & 340{,}000 \\
Master data (combined)    & 58{,}000 \\
Document-flow records     & 85{,}000 \\
Subledger records         & 95{,}000 \\
\bottomrule
\end{tabular}
\end{table}

\paragraph{Distribution Sampling.}
Micro-benchmarks for individual samplers:

\begin{center}
\small
\begin{tabular}{lr}
\toprule
\textbf{Sampler} & \textbf{Latency} \\
\midrule
BenfordSampler          & $<1\,\mu$s / sample \\
Copula (Gaussian, 3-D)  & 2--5\,$\mu$s / sample \\
LogNormalMixture (3-comp.) & $<500$\,ns / sample \\
\bottomrule
\end{tabular}
\end{center}

\paragraph{Output Sink Performance.}

\begin{center}
\small
\begin{tabular}{lr}
\toprule
\textbf{Format} & \textbf{Throughput (records/sec)} \\
\midrule
CSV (streaming)   & 210{,}000 \\
JSON (streaming)  & 105{,}000 \\
Parquet (columnar) & 520{,}000 \\
\bottomrule
\end{tabular}
\end{center}

\paragraph{Scaling.}
Throughput scales near-linearly with core count up to 8 cores (due to the
per-generator seed isolation described in \Cref{sec:arch_determinism}),
after which contention on shared data structures (balance tracker,
document-chain manager) introduces sublinear scaling.  At 16 cores, the
Manufacturing (Medium) configuration generates a full 24-month dataset
(385K journal entries, 1.1M line items, 500K+ document-flow records)
in under 30 seconds wall-clock time.

\paragraph{Memory Efficiency.}
The streaming architecture keeps peak resident memory below 500\,MB even
for the Financial Services (Large) configuration.  The graceful degradation
system (\Cref{sec:arch_resources}) was not triggered during any benchmark
run on the 64\,GB test machine.

% ─────────────────────────────────────────────────────────────────────
\subsection{Privacy-Utility Evaluation}
\label{sec:exp_privacy}

We evaluate the fingerprint module by extracting a fingerprint from the
Manufacturing (Medium) dataset at each privacy level, synthesizing a new
dataset from each fingerprint, and comparing statistical properties.

\begin{table}[t]
\centering
\caption{Fingerprint privacy-utility tradeoff.  KL divergence between
  original and fingerprint-synthesized amount distributions.}
\label{tab:privacy_utility}
\small
\begin{tabular}{lccc}
\toprule
\textbf{Privacy Level} & $\varepsilon$ & \textbf{KL Divergence}
  & \textbf{Benford MAD} \\
\midrule
Minimal  & 5.0 & 0.003 & 0.0026 \\
Standard & 1.0 & 0.008 & 0.0039 \\
High     & 0.5 & 0.019 & 0.0052 \\
Maximum  & 0.1 & 0.071 & 0.0098 \\
\bottomrule
\end{tabular}
\end{table}

Even at the ``High'' privacy level ($\varepsilon = 0.5$), the
KL divergence remains below 0.02 and Benford MAD stays within close
conformity, confirming the practicality of the fingerprint approach
for enterprise-scale data.
